\chapter{Conclusioni}

Il presente elaborato ha esaminato la modernizzazione del primo stadio del \textit{Drowning Early Warning System}, dedicato alla rilevazione e alla segmentazione dei bagnanti. L'obiettivo non era sostituire l'intera architettura di DEWS, bensì valutare se un modello della famiglia YOLO potesse fornire rilevazioni meno dipendenti dalla configurazione manuale, sufficientemente accurate e compatibili con un'elaborazione in tempo reale, conservando le informazioni necessarie al tracciamento e alla successiva analisi del comportamento.

Il lavoro ha riguardato l'intero percorso sperimentale richiesto da tale valutazione: dalla realizzazione di un \textit{dataset} specifico per l'ambiente acquatico al \textit{fine-tuning} di modelli per la \textit{detection} e la segmentazione di istanza, fino all'analisi del loro possibile inserimento nella \textit{pipeline} esistente. Le conclusioni che seguono valutano il raggiungimento degli obiettivi iniziali e delimitano la validità dei risultati ottenuti.

\section{Valutazione rispetto agli obiettivi iniziali}

Il primo obiettivo, relativo alla realizzazione di un \textit{dataset} rappresentativo delle difficoltà visive dell'ambiente acquatico, può essere considerato raggiunto entro il perimetro sperimentale definito. La raccolta comprende infatti scene realistiche con riflessi, schizzi, occlusioni, oggetti e soggetti parzialmente immersi e permette di addestrare e confrontare in modo controllato modelli di \textit{detection} e segmentazione.

È stato raggiunto anche l'obiettivo di valutare la sostituzione del rilevatore originale mediante YOLO. I risultati sul \textit{test set} mostrano che l'adattamento al dominio consente di individuare i bagnanti con metriche elevate e con tempi di inferenza contenuti sull'\textit{hardware} di prova. Il requisito temporale risulta quindi plausibile per il singolo modello: persino la configurazione più lenta fra quelle esaminate mantiene un margine ampio rispetto alla frequenza di acquisizione prevista per ciascuna telecamera.

L'obiettivo di compatibilità con le fasi successive è stato verificato sul piano informativo e architetturale. YOLO-SEG fornisce direttamente una maschera distinta per ogni persona, mentre YOLO \textit{Detection} richiede uno stadio aggiuntivo per ricostruire la \textit{silhouette}. In entrambi i casi è possibile preservare l'interfaccia verso il tracciamento e il calcolo dei descrittori. Tale risultato non coincide tuttavia con la validazione completa del sistema: non sono state misurate la latenza \textit{end-to-end}, la stabilità delle tracce né l'accuratezza finale nel riconoscimento delle situazioni di crisi. Il lavoro dimostra quindi l'idoneità del nuovo modulo come candidato alla sostituzione del rilevatore, non ancora l'efficacia operativa di un sistema completo di allarme.

\section{Limiti dell'approccio proposto}

Il limite principale riguarda la rappresentatività dei dati. Tutte le registrazioni provengono da un solo impianto e da una sola giornata di acquisizione. Le inquadrature del \textit{test set} non compaiono durante l'addestramento, ma condividono con gli altri insiemi la piscina, le condizioni generali di illuminazione e il contesto architettonico. Le prestazioni ottenute misurano pertanto la generalizzazione verso punti di vista non osservati, ma non permettono di stabilire come il modello si comporti in piscine differenti o in presenza di variazioni ambientali più ampie.

La qualità delle annotazioni non è inoltre uniforme. Il procedimento semiautomatico ha consentito di privilegiare l'estensione della raccolta, ma alcune maschere conservano contorni imprecisi o regioni parziali. A ciò si aggiunge la correlazione fra fotogrammi consecutivi, che rende il \textit{dataset} numericamente esteso ma in parte ridondante dal punto di vista dell'informazione contenuta.

Anche la campagna di addestramento rimane circoscritta. Non tutte le combinazioni fra taglia, compito e configurazione sono state sperimentate nelle medesime condizioni, alcuni addestramenti non si sono conclusi secondo il criterio previsto e la valutazione ha considerato una sola risoluzione di ingresso, pari a 640 \textit{pixel}. Le metriche riportate descrivono inoltre la qualità delle predizioni sui singoli fotogrammi; non misurano direttamente la continuità temporale delle rilevazioni, la durata degli errori o il loro effetto sui descrittori comportamentali.

Infine, i tempi di inferenza misurati riguardano soltanto l'esecuzione del modello e non rappresentano la latenza complessiva della \textit{pipeline}. Non è quindi possibile ricavare dai soli valori misurati il numero di flussi video gestibili contemporaneamente né confrontare in modo definitivo YOLO-SEG con la soluzione ibrida basata su \textit{detection} e segmentazione locale.

\section{Possibili sviluppi futuri}

Un primo sviluppo consiste nell'estendere la raccolta a più piscine, includendo ambienti interni ed esterni, differenti profondità, illuminazioni, caratteristiche dell'acqua e dispositivi di acquisizione. La separazione dei dati dovrebbe essere mantenuta a livello di impianto e di sessione, così da ottenere una misura più rigorosa della capacità di generalizzazione. Una revisione mirata delle annotazioni più incerte e una selezione meno ridondante dei fotogrammi potrebbero inoltre migliorare il rapporto fra qualità e dimensione del \textit{dataset}.

Sul piano sperimentale, sarà utile completare il confronto fra taglie e configurazioni, valutare risoluzioni di ingresso differenti e analizzare gli errori in funzione dell'affollamento, del grado di immersione e delle occlusioni. L'osservazione di lunghe sequenze continue, oltre che di singoli fotogrammi, permetterebbe di misurare la persistenza dei falsi positivi e dei falsi negativi e di regolare la soglia di confidenza in funzione della stabilità delle tracce.

Il passaggio successivo più rilevante rimane l'integrazione effettiva di YOLO26-S-SEG nella \textit{pipeline}. Tale implementazione dovrà verificare il ripristino delle maschere nelle coordinate originali, l'associazione temporale delle istanze e la coerenza dei descrittori rispetto al sistema di riferimento. Un \textit{benchmark end-to-end}, eseguito sull'\textit{hardware} destinato all'impiego e con più flussi video contemporanei, consentirà di misurare latenza, frequenza effettiva di elaborazione e comportamento nei casi più gravosi.

In parallelo potrà essere confrontata la soluzione diretta basata su YOLO-SEG con il metodo ibrido, nel quale YOLO \textit{Detection} individua le persone e la segmentazione viene applicata soltanto all'interno delle regioni rilevate. La definizione di un'area di interesse coincidente con la vasca potrà ridurre ulteriormente le elaborazioni non necessarie e le rilevazioni esterne alla scena utile. Il confronto dovrà considerare non soltanto il tempo medio, ma anche la qualità delle maschere e l'effetto complessivo sul riconoscimento del comportamento.

Nel complesso, i risultati ottenuti indicano che l'introduzione di YOLO costituisce una direzione concreta per semplificare e aggiornare il primo stadio di DEWS. Il lavoro fornisce il \textit{dataset}, i modelli adattati e l'analisi necessaria a selezionare una configurazione iniziale; la validazione in nuovi ambienti e l'esecuzione dell'intera \textit{pipeline} rappresentano i passaggi necessari per trasformare tale risultato sperimentale in un sistema valutabile in condizioni operative.
